\subsection{Measurement Integrity}
\label{sec:integrity-theory}
The competition data also contain integrity events: frames that never
arrive and frames whose values are corrupted. Write the received report of
PMU $p$ as
\begin{equation}
 y_{p,k}=M_{p,k}\big(h_p(x_k)+v_{p,k}+\beta_{p,k}\big),
 \label{eq:integrity-observation}
\end{equation}
where $M_{p,k}\in\{0,1\}$ is availability and $\beta_{p,k}$ a corruption
that does not enter the network equations. Availability is observed through
the frame flag, so a missing PMU is identified exactly. What it removes is
information, and the same quantity governs corruption.

\begin{proposition}[Integrity isolation]\label{prop:integrity}
Let the noise be Gaussian and independent across PMUs, and let the
corruption of PMU $p$ be unrestricted over a window. The minimum Gaussian divergence from physical source $i$ from ``no
intervention, PMU $p$ corrupted'' equals
$\Lambda_{i0}^{(-p)}$, the detection information of $i$ computed without the
channels of PMU $p$. This is also the information for $i$ when PMU $p$ is
missing.
\end{proposition}
\emph{Proof.} Minimizing over an unrestricted $\beta_p$ projects the signature onto
the orthogonal complement of the channels of PMU $p$. With block-diagonal
noise covariance this projection deletes those channels, and the remaining
quadratic form is the information of the other PMUs. \hfill$\square$

For a common joint Gaussian covariance $\Sigma$ and a free shift on the
corrupted coordinates, minimizing the quadratic separation over that
shift gives $r_c^\top\Sigma_{cc}^{-1}r_c$, where $c$ denotes the remaining
coordinates. This is the marginal covariance of the clean observations.
Equivalently, it is the Schur complement of the corrupted block in the
joint \emph{precision} matrix. Conditioning on the corrupted values as
though they were clean would answer a different question. The
proposition explains why integrity cannot be recognized from a single PMU
alone: a corruption becomes distinguishable only through relations that
every physical event of the network obeys. Two such relations are used.
\emph{Sequence parity.} For balanced interventions (load, line,
generation, three-phase fault) the balanced network produces no negative-
or zero-sequence voltage or current. These sequence components are
therefore ideal-model parity relations in the sense of analytical redundancy
\cite{chow1984analytical}: they vanish in this balanced observation model and sensitive to corruption of a single phase. A step
$\epsilon$ in one phase magnitude gives $|X_2|=|X_0|=\epsilon|X|/3$. The
corrupted phase is the one whose phasor is most aligned with $X_0$.
Under independent per-phase noise, $q_k=(|X_2|^2+|X_0|^2)/\sigma^2$ is
$\mathrm{Gamma}(2,1)$-distributed, which fixes the threshold for a given
false-alarm probability. \emph{Network coupling.} A modeled single bolted fault that changes the positive-sequence voltage of PMU $p$ by
$\Delta_p$ changes the other PMUs in proportion to the network impedance
matrix $Z=Y^{-1}$. Here $Y$ includes constant-impedance loads and generator
transient reactances \cite{pai1989energy}, and a bolted fault at bus $f$
gives $\Delta_j/\Delta_p=|Z_{jf}/Z_{pf}|$. Hence
\begin{equation}
 \max_{j\neq p}\Delta_j\ \ge\ c_p\,\Delta_p,\qquad
 c_p=\min_{f\in\mathcal F_p}\max_{j\neq p}\Big|\frac{Z_{jf}}{Z_{pf}}\Big|,
 \label{eq:coupling-bound}
\end{equation}
with $\mathcal F_p$ the buses at which $p$ is the most affected PMU. For the
IEEE 39-bus network $c_p$ ranges from \IntCouplingMin{} (PMU~29) to
\IntCouplingMax{}. A balanced excursion that violates
\eqref{eq:coupling-bound} by a wide margin is inconsistent with the modeled set of single bolted faults.
It does not rule out other physical disturbances or report-window effects.
Both relations are classical bad-data reasoning \cite{abur2004state},
applied here to report-rate phasors without a state estimator.
